\documentclass[12pt]{article}
\usepackage{amsmath}
\usepackage{amsfonts}
\usepackage{float}
\usepackage{fancyhdr}
\usepackage{graphicx}
\usepackage{setspace}
\usepackage[colorlinks=true,linkcolor=blue, citecolor=red]{hyperref}
\usepackage{url}
\usepackage[top=.75in, left=.75in, right=.75in, bottom=1in]{geometry}

% For algorithm
\usepackage{algorithm}
\usepackage{algpseudocode}

% ============ CODE ============
\usepackage{listings}
\usepackage{xcolor}
\definecolor{codegreen}{rgb}{0,0.6,0}
\definecolor{codegray}{rgb}{0.5,0.5,0.5}
\definecolor{codepurple}{rgb}{0.58,0,0.82}
\definecolor{backcolour}{rgb}{0.95,0.95,0.92}

% Styling for the code.
\lstdefinestyle{mystyle}{
    backgroundcolor=\color{backcolour},   
    commentstyle=\color{codegreen},
    keywordstyle=\color{magenta},
    numberstyle=\tiny\color{codegray},
    stringstyle=\color{codepurple},
    basicstyle=\ttfamily\footnotesize,
    breakatwhitespace=false,         
    breaklines=true,                 
    captionpos=b,                    
    keepspaces=true,                 
    numbers=left,                    
    numbersep=5pt,                  
    showspaces=false,                
    showstringspaces=false,
    showtabs=false,                  
    tabsize=2
}
\lstset{style=mystyle}

% Table packages and Unicode / Font support
\usepackage{array}
\usepackage{longtable}
\usepackage{enumitem}
\usepackage{booktabs}
\usepackage[utf8]{inputenc}
\DeclareUnicodeCharacter{2264}{\ensuremath{\le}}
\DeclareUnicodeCharacter{2194}{\ensuremath{\leftrightarrow}}
\usepackage[T1]{fontenc}
\usepackage{lmodern}
\graphicspath{{../../../../templates/}{../../prototype/}{figures/}{./}}
\setlength{\emergencystretch}{3em}
\setlist[enumerate]{label=\arabic*.,leftmargin=*}
\setlist[itemize]{leftmargin=*}

% ============ HEADER AND FOOTER ============
\setlength{\headheight}{29.43912pt}

\pagestyle{fancy}
\fancyfoot{}
\fancyfoot[R]{Page \thepage}

\setlength{\parindent}{0pt}
\renewcommand{\baselinestretch}{1.5}

\newcommand{\coursename}{Introduction of Software Engineering}
\newcommand{\reportname}{Studitopia}
\newcommand{\reporttitle}{\href{https://github.com/RoundFloor/Studitopia.git}{Studitopia}}

\newcommand{\studentname}{%
Nguyen Thanh An (24125024)\\
Nguyen Chi Bao (24125025)\\
Tran Nguyen Khanh Binh (24125026)\\
Tran Canh Anh Tuan (24125048)\\
Nguyen Le Quoc Huy (24125100)
}
\newcommand{\teachername}{Ho Tuan Thanh, Nguyen Thi Minh Tuyen}

% Header
\lhead{\reporttitle}
\rhead{
University of Science - VNUHCM\\
\coursename
}

% Footer
\newcommand{\leftfooter}{}
\lfoot{\leftfooter}

% ============ DOCUMENT ============
\begin{document}

\pagenumbering{roman}
\begin{titlepage}
\newcommand{\HRule}{\rule{\linewidth}{0.5mm}}
\centering

\textsc{\LARGE VNU-HCM}\\[0.5cm]
\textsc{\Large University of Science}\\[0.5cm]
\textsc{\large Faculty of Information Technology}\\[0.5cm]
\textsc{Computer Science}\\[0.5cm]

\HRule \\[0.4cm]
{ 
\huge{\bfseries{\reporttitle}}\\[0.5cm]
\large{\bfseries{UC-03 --- Browse and Search the University Forum}}
}\\[0.4cm]
\HRule \\[0.5cm]

\textbf{\large Subject: \coursename}\\[0.5cm]

\begin{minipage}[t]{0.55\textwidth}
\begin{flushleft} \large
\emph{Students:}\\
\studentname
\end{flushleft}
\end{minipage}
~
\begin{minipage}[t]{0.4\textwidth}
\begin{flushright} \large
\emph{Supervising Teachers:} \\
\teachername
\end{flushright}
\end{minipage}\\[1cm]

{\large \today}\\[1cm]

\includegraphics[scale=.20]{img/hcmus-logo.png}\\[1cm] 

\vfill
\end{titlepage}	

\tableofcontents
\pagebreak

\pagenumbering{arabic}
\setcounter{page}{1}

\hypertarget{sec:uc-03-browse-and-search-the-university-forum}{}
\addcontentsline{toc}{section}{UC-03 --- Browse and Search the University Forum}
\section*{UC-03 --- Browse and Search the University Forum}
\par\medskip

\textbf{Version:} \texttt{2.0}
\par\medskip

\textbf{Performed by:} [Nguyen Chi Bao] | \textbf{Edited by:} [Nguyen Le Quoc Huy]
\par\medskip

\begingroup
\small
\renewcommand{\baselinestretch}{1.0}\selectfont
\renewcommand{\arraystretch}{1.15}
\setlength{\tabcolsep}{6pt}
\begin{longtable}{|>{\raggedright\arraybackslash}p{0.09\textwidth}|>{\raggedright\arraybackslash}p{0.84\textwidth}|}
\hline
\textbf{Field} & \textbf{Value} \\ \hline
\endhead
Use Case ID & UC-03 \\ \hline
Name & Browse and Search the University Forum \\ \hline
Primary Actor & University Student \\ \hline
Supporting Actors & None \\ \hline
Feature Group & Forum (Primary) \\ \hline
Priority & High \\ \hline
Related User Stories & US-11, US-12, US-13, US-17, US-23 \\ \hline
Related Requirements & US-11, US-12, US-13, US-17, US-23; FR-11.1 to FR-11.4, FR-12.1 to FR-12.6, FR-13.1 to FR-13.5, FR-17.1 to FR-17.5, FR-23.1 to FR-23.5; NFR-PERF-2, NFR-PERF-3, NFR-PERF-5, NFR-SEC-4, NFR-SEC-6, NFR-PRIV-7, NFR-USE-3 \\ \hline
Source Documents & \texttt{proposals/use\_cases\_model/UC-03-browse-search-forum.md}; \texttt{proposals/use\_cases\_description.md} UC-03; \texttt{proposals/use\_cases\_model.md}; \texttt{proposals/use\_case\_diagrams.md}; \texttt{PA3-Group05/use\_case.md} §3; \texttt{proposals/overview.md}; \texttt{proposals/functional\_requirement.md}; \texttt{proposals/non\_functional\_requirement.md}; \texttt{proposals/userstory.md}; \texttt{proposals/feature\_flow.md}; \texttt{proposals/screen\_flow.md}; \texttt{proposals/API.md}; \texttt{proposals/response\_code.md}; \texttt{proposals/list\_post.md}; \texttt{proposals/post\_management.md} \\ \hline
\end{longtable}
\endgroup
\par\medskip

\hypertarget{sec:part-a-use-case-model}{}
\addcontentsline{toc}{subsection}{Part A --- Use Case Model}
\subsection*{Part A --- Use Case Model}
\par\medskip

\hypertarget{sec:a1-focused-use-case-diagram}{}
\addcontentsline{toc}{subsubsection}{A1. Focused Use Case Diagram}
\subsubsection*{A1. Focused Use Case Diagram}
\par\medskip

\begin{center}
\includegraphics[width=0.85\textwidth,height=0.75\textheight,keepaspectratio]{"figures/UC-03-browse-search-forum_1.png"}
\end{center}
\par\medskip

\textbf{Diagram notation.} Person shape = human actor; the system boundary contains the use cases; blue stadium = selected base use case UC-03; purple stadium = included use case I-04; green stadium = extending use cases E-03 and E-04; solid line = association; dashed arrow = include or extend. The diagram is copied verbatim from the focused UC-03 model.
\par\medskip

\hypertarget{sec:a2-actor-register}{}
\addcontentsline{toc}{subsubsection}{A2. Actor Register}
\subsubsection*{A2. Actor Register}
\par\medskip

\begingroup
\small
\renewcommand{\baselinestretch}{1.0}\selectfont
\renewcommand{\arraystretch}{1.15}
\setlength{\tabcolsep}{6pt}
\begin{longtable}{|>{\raggedright\arraybackslash}p{0.15\textwidth}|>{\raggedright\arraybackslash}p{0.12\textwidth}|>{\raggedright\arraybackslash}p{0.65\textwidth}|}
\hline
\textbf{Actor} & \textbf{Type} & \textbf{Description} \\ \hline
\endhead
University Student & Primary human actor & A registered student who browses university, major, and course forums, searches academic content, and opens posts to read details. \\ \hline
New University Student & Specialized human actor & Inherits the University Student association to UC-03 through actor generalization. \\ \hline
Registered University Student & Specialized human actor & Inherits the University Student association to UC-03 through actor generalization. \\ \hline
Post Author & Specialized human actor & Inherits the University Student association to UC-03 through actor generalization. \\ \hline
Room Owner or Manager & Specialized human actor & Inherits the University Student association to UC-03 through actor generalization. \\ \hline
\end{longtable}
\endgroup
\par\medskip

\hypertarget{sec:a3-relationships}{}
\addcontentsline{toc}{subsubsection}{A3. Relationships}
\subsubsection*{A3. Relationships}
\par\medskip

\hypertarget{sec:a3-1-associations}{}
\paragraph*{A3.1 Associations}
\par\medskip

\begingroup
\small
\renewcommand{\baselinestretch}{1.0}\selectfont
\renewcommand{\arraystretch}{1.15}
\setlength{\tabcolsep}{6pt}
\begin{longtable}{|>{\raggedright\arraybackslash}p{0.14\textwidth}|>{\raggedright\arraybackslash}p{0.34\textwidth}|>{\raggedright\arraybackslash}p{0.10\textwidth}|>{\raggedright\arraybackslash}p{0.33\textwidth}|}
\hline
\textbf{Actor} & \textbf{Use Case} & \textbf{Type} & \textbf{Description} \\ \hline
\endhead
University Student & UC-03 Browse and Search the University Forum & Primary & Browses and searches academic discussions. \\ \hline
\end{longtable}
\endgroup
\par\medskip

\hypertarget{sec:a3-2-generalization-and-specialization}{}
\paragraph*{A3.2 Generalization and Specialization}
\par\medskip

\begingroup
\small
\renewcommand{\baselinestretch}{1.0}\selectfont
\renewcommand{\arraystretch}{1.15}
\setlength{\tabcolsep}{6pt}
\begin{longtable}{|>{\raggedright\arraybackslash}p{0.42\textwidth}|>{\raggedright\arraybackslash}p{0.10\textwidth}|>{\raggedright\arraybackslash}p{0.10\textwidth}|>{\raggedright\arraybackslash}p{0.29\textwidth}|}
\hline
\textbf{Child} & \textbf{Parent} & \textbf{Type} & \textbf{Inherited Behavior or Association} \\ \hline
\endhead
New University Student; Registered University Student; Post Author; Room Owner or Manager & University Student & Actor generalization & Each specialized student actor inherits the UC-03 association. \\ \hline
\end{longtable}
\endgroup
\par\medskip

\hypertarget{sec:a3-3-include}{}
\paragraph*{A3.3 Include}
\par\medskip

\begingroup
\small
\renewcommand{\baselinestretch}{1.0}\selectfont
\renewcommand{\arraystretch}{1.15}
\setlength{\tabcolsep}{6pt}
\begin{longtable}{|>{\raggedright\arraybackslash}p{0.32\textwidth}|>{\raggedright\arraybackslash}p{0.18\textwidth}|>{\raggedright\arraybackslash}p{0.42\textwidth}|}
\hline
\textbf{Base Use Case} & \textbf{Included Use Case} & \textbf{Meaning} \\ \hline
\endhead
UC-03 Browse and Search the University Forum & I-04 Load Forum Post List & Load the relevant university, major, or course post list. \\ \hline
\end{longtable}
\endgroup
\par\medskip

\hypertarget{sec:a3-4-extend}{}
\paragraph*{A3.4 Extend}
\par\medskip

\begingroup
\small
\renewcommand{\baselinestretch}{1.0}\selectfont
\renewcommand{\arraystretch}{1.15}
\setlength{\tabcolsep}{6pt}
\begin{longtable}{|>{\raggedright\arraybackslash}p{0.10\textwidth}|>{\raggedright\arraybackslash}p{0.20\textwidth}|>{\raggedright\arraybackslash}p{0.26\textwidth}|>{\raggedright\arraybackslash}p{0.36\textwidth}|}
\hline
\textbf{Extending Use Case} & \textbf{Base Use Case} & \textbf{Condition} & \textbf{Added Behavior} \\ \hline
\endhead
E-03 Search Forum & UC-03 Browse and Search the University Forum & The student performs a keyword or filtered search. & Add keyword, tag, author, and academic-context search with ranking. \\ \hline
E-04 View Post Detail & UC-03 Browse and Search the University Forum & The student opens a post from the list or search results. & Add full post detail, comments, answers, true response, attachments, and tools. \\ \hline
\end{longtable}
\endgroup
\par\medskip

\hypertarget{sec:a4-model-notes}{}
\addcontentsline{toc}{subsubsection}{A4. Model Notes}
\subsubsection*{A4. Model Notes}
\par\medskip

\begin{enumerate}
\item The base browse path always includes I-04 because every browse or search experience is based on a forum post list.
\item Search and post detail are optional extensions because a student may browse without searching or opening a post.
\item UC-03 is part of the primary forum model; all specialized student actors inherit its association through University Student.
\item The focused diagram shows only the base actor University Student; the inherited associations of specialized student actors are documented in A3.2 rather than drawn as separate links.
\item Forum content is organized by University, then Major, then Course; each level must show only relevant posts and materials.
\item E-03 and E-04 correspond to the optional search and detail steps in the basic flow, not to separate primary-actor use cases.
\end{enumerate}
\par\medskip

\hypertarget{sec:part-b-use-case-specification}{}
\addcontentsline{toc}{subsection}{Part B --- Use Case Specification}
\subsection*{Part B --- Use Case Specification}
\par\medskip

\hypertarget{sec:b1-brief-description}{}
\addcontentsline{toc}{subsubsection}{B1. Brief Description}
\subsubsection*{B1. Brief Description}
\par\medskip

A student discovers relevant university, major, and course discussions, searches for academic content, and opens a post to read its full details. The forum is the primary collaboration feature of Studitopia. The trigger is opening \textbf{Forum} or submitting a forum search; the successful outcome is a relevant post list, ranked search results, a clear empty state, or a readable post detail view.
\par\medskip

\hypertarget{sec:b2-preconditions}{}
\addcontentsline{toc}{subsubsection}{B2. Preconditions}
\subsubsection*{B2. Preconditions}
\par\medskip

\begin{enumerate}
\item The student is authenticated with an active session.
\item University, major, course, and forum data are available.
\item The student has permission to view the selected forum content.
\end{enumerate}
\par\medskip

\hypertarget{sec:b3-postconditions}{}
\addcontentsline{toc}{subsubsection}{B3. Postconditions}
\subsubsection*{B3. Postconditions}
\par\medskip

\begin{enumerate}
\item \textbf{Success:} The selected forum level and a relevant post list are displayed.
\item \textbf{Success:} Search and filters return ranked results, or a clear empty state is shown when nothing matches.
\item \textbf{Success:} The student can open a post and read its full content, metadata, discussion, attachments, and permitted actions.
\item \textbf{Failure:} The system displays an error or permission message without exposing restricted content.
\end{enumerate}
\par\medskip

\hypertarget{sec:b4-main-success-scenario-basic-flow}{}
\addcontentsline{toc}{subsubsection}{B4. Main Success Scenario (Basic Flow)}
\subsubsection*{B4. Main Success Scenario (Basic Flow)}
\par\medskip

\begin{enumerate}
\item The student selects \textbf{Forum} from the main navigation.
\item The system opens the student's default university forum and displays the Forum entry, forum hierarchy, and post list.
\item The student selects a university; the system displays the majors offered by that university.
\item The student selects a major; the system displays the courses associated with that major.
\item The student selects a course; the system displays the course forum and its posts.
\item The system shows each post summary with title, author, post type, timestamp, vote score, comment count, and solved status.
\item The student optionally applies filters such as My Post, area, university, major, course, tags, and post type, and selects a sorting order such as newest, top trending, most voted, or most discussed.
\item The student optionally enters a keyword and searches by tag, title, content, author, university, major, or course.
\item The system ranks matching results by relevance and engagement and displays the filtered result set.
\item The student opens a post from the list or search results.
\item The system displays the full post content, author, post type, forum context, timestamp, vote score, tags, comments, answers, true response, attachments, and functional tools when present.
\item The system displays only the actions permitted for the student, such as vote, comment, save, report, or mark true response.
\end{enumerate}
\par\medskip

\hypertarget{sec:b5-alternative-flows}{}
\addcontentsline{toc}{subsubsection}{B5. Alternative Flows}
\subsubsection*{B5. Alternative Flows}
\par\medskip

Unless otherwise stated, an alternative flow resumes at the next step of the basic flow.
\par\medskip

\hypertarget{sec:a1-browse-another-university}{}
\paragraph*{A1. Browse another university}
\par\medskip

At basic flow step 3, the student selects another university forum when permitted. The system applies that university as the active browsing context and shows only content relevant to that context.
\par\medskip

\hypertarget{sec:a2-change-forum-context}{}
\paragraph*{A2. Change forum context}
\par\medskip

The student switches directly between university, major, and course levels. The system updates the post list to match the selected level.
\par\medskip

\hypertarget{sec:a3-no-posts-match}{}
\paragraph*{A3. No posts match}
\par\medskip

At basic flow step 9 or 10, no posts match the filters or search query. The system displays an empty state and offers to clear filters or change the search terms.
\par\medskip

\hypertarget{sec:a4-search-without-a-keyword}{}
\paragraph*{A4. Search without a keyword}
\par\medskip

The student applies only filters. The system returns results based on the selected forum context and filters.
\par\medskip

\hypertarget{sec:a5-return-to-a-previous-context}{}
\paragraph*{A5. Return to a previous context}
\par\medskip

The student navigates back; the system restores the previous forum level and result state where supported.
\par\medskip

\hypertarget{sec:b6-exception-flows}{}
\addcontentsline{toc}{subsubsection}{B6. Exception Flows}
\subsubsection*{B6. Exception Flows}
\par\medskip

\hypertarget{sec:e1-forum-content-cannot-be-loaded}{}
\paragraph*{E1. Forum content cannot be loaded}
\par\medskip

The system displays an error state and a retry option without showing partial or stale restricted content.
\par\medskip

\hypertarget{sec:e2-selected-post-is-unavailable-or-deleted}{}
\paragraph*{E2. Selected post is unavailable or deleted}
\par\medskip

The system informs the student that the post is no longer available and returns them to the post list.
\par\medskip

\hypertarget{sec:e3-access-is-not-permitted}{}
\paragraph*{E3. Access is not permitted}
\par\medskip

The system denies access to private or restricted content and explains that the student lacks permission.
\par\medskip

\hypertarget{sec:e4-invalid-search-or-filter-combination}{}
\paragraph*{E4. Invalid search or filter combination}
\par\medskip

The system displays validation feedback and preserves the usable parts of the query.
\par\medskip

\hypertarget{sec:b7-business-rules}{}
\addcontentsline{toc}{subsubsection}{B7. Business Rules}
\subsubsection*{B7. Business Rules}
\par\medskip

\begin{enumerate}
\item Forum content is organized by University, then Major, then Course.
\item After login, the student's own university is the default forum context.
\item Each forum level must display only posts and materials relevant to that level.
\item Post summaries must include title, author, post type, timestamp, vote score, comment count, and solved status.
\item Search supports tags, title, content, author, university, major, and course.
\item Search results are ranked by relevance and engagement and can be filtered by post type, solved status, and time range.
\item Post detail must show the full content, metadata, discussion, marked true response, attachments, and functional tools when present.
\item Available post actions are determined by the student's permissions.
\item Privacy and forum-access rules apply consistently to post lists, post details, and search results.
\end{enumerate}
\par\medskip

\hypertarget{sec:b8-special-requirements}{}
\addcontentsline{toc}{subsubsection}{B8. Special Requirements}
\subsubsection*{B8. Special Requirements}
\par\medskip

\begingroup
\small
\renewcommand{\baselinestretch}{1.0}\selectfont
\renewcommand{\arraystretch}{1.15}
\setlength{\tabcolsep}{6pt}
\begin{longtable}{|>{\raggedright\arraybackslash}p{0.10\textwidth}|>{\raggedright\arraybackslash}p{0.83\textwidth}|}
\hline
\textbf{ID} & \textbf{Requirement} \\ \hline
\endhead
NFR-PERF-2 & A forum post list must load within 2 seconds under normal operating conditions. \\ \hline
NFR-PERF-3 & A post detail page and its comments must display within 2 seconds. \\ \hline
NFR-PERF-5 & Forum search results must return within 2 seconds for common queries. \\ \hline
NFR-SEC-4 & Private forum content and private-group content must be restricted to authorized users only. \\ \hline
NFR-SEC-6 & Server-side authorization checks must prevent unauthorized access. \\ \hline
NFR-PRIV-7 & Privacy settings must apply consistently to forum posts, comments, member lists, and search results. \\ \hline
NFR-USE-3 & Navigation and interaction patterns must remain consistent across university, major, and course forums. \\ \hline
\end{longtable}
\endgroup
\par\medskip

\hypertarget{sec:b9-extension-points}{}
\addcontentsline{toc}{subsubsection}{B9. Extension Points}
\subsubsection*{B9. Extension Points}
\par\medskip

\begingroup
\small
\renewcommand{\baselinestretch}{1.0}\selectfont
\renewcommand{\arraystretch}{1.15}
\setlength{\tabcolsep}{6pt}
\begin{longtable}{|>{\raggedright\arraybackslash}p{0.33\textwidth}|>{\raggedright\arraybackslash}p{0.60\textwidth}|}
\hline
\textbf{Extension Point} & \textbf{Location in Flow of Events} \\ \hline
\endhead
Change browsing context & Basic flow steps 3-5; alternatives A1-A2 \\ \hline
Apply filters & Basic flow step 7; alternative A4 \\ \hline
Search forum & Basic flow step 8; alternative A3 and exception E4 \\ \hline
View post detail & Basic flow step 10; exception E2 \\ \hline
Empty-result recovery & Basic flow step 9; alternative A3 \\ \hline
Restricted-content handling & Access check; exception E3 \\ \hline
\end{longtable}
\endgroup
\par\medskip

\hypertarget{sec:b10-interface-and-data-notes}{}
\addcontentsline{toc}{subsubsection}{B10. Interface and Data Notes}
\subsubsection*{B10. Interface and Data Notes}
\par\medskip

\begin{itemize}
\item \textbf{Screen flow.} \texttt{Forum → University → Major → Course → Post}; Home has the \textbf{Forum} entry; \textbf{Trending $\rightarrow$ Post} and \textbf{Trending $\rightarrow$ Related $\rightarrow$ Post} are also shown in \texttt{screen\_flow.md}. \texttt{feature\_flow.md} adds \texttt{Forum → Recommended → List Post}, \texttt{Forum → Search}, \texttt{List Post → Filter → My Post / Area / University / Major / Course / Tags}, \texttt{List Post → Post}, and \texttt{Post → Comment / Mark True Response / Upvote / Downvote / Save / Report / Moderate}. \texttt{list\_post.md} proposes \texttt{ListPostPage} / \texttt{ForumPostListPage}, \texttt{/forum}, \texttt{/forum/:universityId/:majorId/:courseId}, and \texttt{/forum/posts/:postId} as recommendations, not source-confirmed routes.
\item \textbf{Post-card summary fields.} Required summary fields are title, author, post type, timestamp, vote score, comment count, and solved status. \texttt{API.md} returns \texttt{post\_id}, \texttt{title}, \texttt{post\_type}, \texttt{author}, \texttt{vote\_score}, \texttt{comment\_count}, \texttt{solved}, \texttt{created\_at}, and \texttt{total}. \texttt{list\_post.md} additionally recommends tags, forum context badge, attachment indicator, and functional-tool badge; these are recommendation-layer fields.
\item \textbf{Filter and sort fields.} Filters: My Post, area, university, major, course, tags, post type (\texttt{Discussion}, \texttt{Material}, \texttt{Functional}), solved status (\texttt{All}, \texttt{Solved}, \texttt{Unsolved}), and time range (\texttt{Today}, \texttt{7 days}, \texttt{30 days}, \texttt{All time}). Sorting: newest, top trending, most voted, most discussed; search results are ranked by relevance and engagement.
\item \textbf{Search fields.} Keyword plus tags, title, content, author, university, major, and course under FR-23.1 and FR-23.2. \texttt{GET /forum/posts} declares \texttt{university\_id?}, \texttt{major\_id?}, \texttt{course\_id?}, \texttt{tag?}, \texttt{search?}, \texttt{sort?} = \texttt{'newest' | 'trending' | 'most\_voted' | 'most\_discussed'}, \texttt{limit?}, and \texttt{offset?}.
\item \textbf{Empty/error states.} Empty post/search results show a clear empty state with clear-filters or change-terms recovery; load failure shows an error banner with \textbf{Retry}; unavailable or deleted post shows a post-unavailable state with return-to-list; restricted content shows an access-denied state; invalid filter/search combinations show validation feedback and preserve usable query parts.
\item \textbf{Recommendation rules.} From FR-24.1 to FR-24.5: recommend posts relevant to university, major, courses, and interests; recommend trending posts within the user's academic context; recommend unsolved questions the user may help with; allow dismissing irrelevant recommendations; refresh recommendations when preferences or forum context change. These rules appear in feature/UI sources but are not listed in the UC-03 Related Requirements field.
\item \textbf{List and search API.} \texttt{GET /forum/posts} $\rightarrow$ \texttt{SUCCESS 200} with posts and total; \texttt{BAD\_REQUEST 400}; \texttt{UNAUTHORIZED 401}; \texttt{INTERNAL\_SERVER\_ERROR 500} with \texttt{(DATABASE\_ERROR | UNKNOWN\_ERROR)[]}. Database function: \texttt{studitopia.list\_forum\_posts(puniversity\_id, pmajor\_id, pcourse\_id, ptag, psearch, psort, plimit, offset)} returning \texttt{post\_id}, \texttt{title}, \texttt{post\_type}, \texttt{author\_id}, \texttt{vote\_score}, \texttt{comment\_count}, \texttt{solved}, \texttt{created\_at}; errors \texttt{UNKNOWN\_ERROR: 45000}, \texttt{INVALID\_POST\_TYPE: 48003}, \texttt{FORBIDDEN: 48010}.
\item \textbf{Detail API.} \texttt{GET /forum/posts/:post\_id} $\rightarrow$ \texttt{SUCCESS 200} with \texttt{post}, \texttt{comments}, \texttt{true\_response}, \texttt{user\_vote} (\texttt{-1 | 0 | 1}), and \texttt{saved}; \texttt{NOT\_FOUND 404} with \texttt{POST\_NOT\_FOUND}; \texttt{UNAUTHORIZED 401}; \texttt{INTERNAL\_SERVER\_ERROR 500} with \texttt{(DATABASE\_ERROR | UNKNOWN\_ERROR)[]}. Database function: \texttt{studitopia.get\_forum\_post(ppost\_id)}; errors \texttt{UNKNOWN\_ERROR: 45000}, \texttt{POST\_NOT\_FOUND: 48001}, \texttt{FORBIDDEN: 48010}.
\item \textbf{Permitted-action endpoints on detail.} \texttt{POST /forum/votes} $\rightarrow$ \texttt{200}, \texttt{400 INVALID\_VOTE}, \texttt{409 DUPLICATED\_VOTE}, \texttt{404 POST\_NOT\_FOUND | COMMENT\_NOT\_FOUND}, \texttt{401}, \texttt{500}; \texttt{POST /forum/posts/:post\_id/save} $\rightarrow$ \texttt{200}, \texttt{404 POST\_NOT\_FOUND}, \texttt{401}, \texttt{500}; \texttt{PATCH /forum/posts/:post\_id/true-response} $\rightarrow$ \texttt{200}, \texttt{404 POST\_NOT\_FOUND | COMMENT\_NOT\_FOUND}, \texttt{403 NOT\_POST\_OWNER | FORBIDDEN}, \texttt{401}, \texttt{500}; \texttt{POST /forum/posts/:post\_id/comments} $\rightarrow$ \texttt{201}, \texttt{400 INVALID\_POST\_CONTENT}, \texttt{404 POST\_NOT\_FOUND | COMMENT\_NOT\_FOUND}, \texttt{401}, \texttt{500}; \texttt{POST /forum/reports} $\rightarrow$ \texttt{201}, \texttt{400 INVALID\_REPORT\_CATEGORY}, \texttt{404 POST\_NOT\_FOUND | COMMENT\_NOT\_FOUND}, \texttt{401}, \texttt{500}. Error-code register includes \texttt{48000 UNKNOWN FORUM ERROR}, \texttt{48001 POST\_NOT\_FOUND}, \texttt{48002 COMMENT\_NOT\_FOUND}, \texttt{48003 INVALID\_POST\_TYPE}, \texttt{48004 INVALID\_POST\_CONTENT}, \texttt{48005 INVALID\_VOTE}, \texttt{48006 DUPLICATED\_VOTE}, \texttt{48007 NOT\_POST\_OWNER}, \texttt{48008 ALREADY\_TRUE\_RESPONSE}, \texttt{48009 INVALID\_REPORT\_CATEGORY}, and \texttt{48010 FORBIDDEN}.
\item \textbf{Privacy and permissions.} Private forum and private-group content must be restricted to authorized users; server-side authorization must prevent unauthorized access; privacy settings must apply consistently to forum posts, comments, member lists, and search results. Other university forums are browsable only when permitted. Detail actions such as vote, comment, save, report, and mark true response must be rendered and enforced according to permissions.
\item \textbf{Persistence and visibility.} \texttt{studitopia.list\_forum\_posts} must return only posts visible to the requesting user; \texttt{studitopia.get\_forum\_post} must respect privacy and moderation visibility rules. The student's own university is the default forum context, and saved posts persist across forum context changes under FR-21.4.
\end{itemize}
\par\medskip

\hypertarget{sec:b11-source-alignment-notes}{}
\addcontentsline{toc}{subsubsection}{B11. Source Alignment Notes}
\subsubsection*{B11. Source Alignment Notes}
\par\medskip

\begin{enumerate}
\item The focused model, \texttt{use\_case\_diagrams.md}, and PA3 §3 contain the same UC-03 Mermaid diagram content; this export reproduces the focused model diagram verbatim.
\item \texttt{use\_cases\_description.md} and PA3 §3 agree on UC-03 goal, preconditions, postconditions, basic flow, five alternatives A1-A5, four exceptions E1-E4, nine business rules, special requirements, and extension points.
\item Performance wording differs slightly: PA3 uses "within 2 seconds" while \texttt{non\_functional\_requirement.md} uses "≤ 2 seconds"; the PA3 special-requirements table is reproduced verbatim above.
\item Recommendations are described in \texttt{feature\_flow.md}, \texttt{list\_post.md}, and FR-24.1 to FR-24.5, but US-24/FR-24 are not in the UC-03 Related Requirements field and recommendations do not appear in the UC-03 basic flow; they are treated as interface/discovery notes only.
\item \texttt{list\_post.md} labels routes such as \texttt{/forum}, \texttt{/forum/:universityId/:majorId/:courseId}, and \texttt{/forum/posts/:postId}, plus tags/attachment/tool badges, as recommendations; the proposal documents do not define confirmed screen routes or those summary fields as FR/AC requirements.
\item \texttt{API.md} exposes \texttt{search}, \texttt{tag}, and academic-context query parameters for \texttt{GET /forum/posts}, but not explicit solved-status or time-range parameters, although business rule 6 and FR-23.3 require those filters; no dedicated API is given for university/major/course catalog listing, My Posts, or saved-post lists.
\item A5 "Return to a previous context" has no separate row in the PA3 extension-point table; it is preserved as an alternative flow because it appears in both UC-03 narrative sources.
\item Session-expiry values conflict across proposal sources: FR-04.1 and US-04 use 15 minutes of inactivity while NFR-SEC-7 uses 30 minutes; neither is listed in the UC-03 Related Requirements field, so no expiry value is added to this use case.
\item \texttt{response\_code.md} lists \texttt{UNKNOWN ERROR 45000} in the database table and \texttt{UNKNOWN DATABASE ERROR 45000} in the backend table; \texttt{API.md} reports \texttt{INTERNAL\_SERVER\_ERROR 500} data as \texttt{(DATABASE\_ERROR | UNKNOWN\_ERROR)[]} for both list and detail endpoints, and the exact names are retained as written in each source.
\end{enumerate}
\par\medskip

\par\medskip
\hypertarget{sec:ui-prototype-specimens}{}
\addcontentsline{toc}{subsection}{UI Prototype Specimens}
\subsection*{UI Prototype Specimens}
\par\medskip
\begin{figure}[H]
\centering
\includegraphics[width=0.75\textwidth,height=0.42\textheight,keepaspectratio]{"../../prototype/browse_home_desktop.png"}
\par\vspace{2pt}
{\footnotesize \textbf{UI Prototype:} Browse home - Desktop}
\end{figure}
\par\smallskip
\begin{figure}[H]
\centering
\includegraphics[width=0.75\textwidth,height=0.42\textheight,keepaspectratio]{"../../prototype/browse_computer_science_desktop.png"}
\par\vspace{2pt}
{\footnotesize \textbf{UI Prototype:} Browse: Computer Science - Desktop}
\end{figure}
\par\smallskip
\begin{figure}[H]
\centering
\includegraphics[width=0.75\textwidth,height=0.42\textheight,keepaspectratio]{"../../prototype/forum_feed_default_desktop.png"}
\par\vspace{2pt}
{\footnotesize \textbf{UI Prototype:} Forum Feed - Default - Desktop}
\end{figure}
\par\smallskip

\end{document}
